\hypertarget{appendix-4fbc03d20829}{}\label{referencestart:appendix-4fbc03d20829}
\section{Two's complement}\label{appendix-4fbc03d20829-twos-complement}

\hypertarget{appendix-4fbc03d20829-twos-complement}{}

\textbf{Curriculum depth:} Developed · \textbf{First seen:} C001 · \textbf{Latest development:} C002

\begin{allmightroledefinition}

\subsection{Definition and mechanism}\label{appendix-4fbc03d20829-definition-and-mechanism}

\hypertarget{appendix-4fbc03d20829-definition-and-mechanism}{}

An \(N\)-bit negative \(-x\) is residue \(2^N-x=\sim x+1\); the signed high-bit weight is \(-2^{N-1}\), giving one zero and an asymmetric range.

\end{allmightroledefinition}

For raw unsigned value u in N bits, signed decoding is u when u\textless2ᴺ⁻¹, otherwise u−2ᴺ. Thus \BookInlineCode{0xAD} means 173−256=−83. Negation uses width-limited complement plus one modulo 2ᴺ. Signed interpretation and unsigned residue share bits but have different language arithmetic contracts.

\subsection{Prerequisites}\label{appendix-4fbc03d20829-prerequisites}

\hypertarget{appendix-4fbc03d20829-prerequisites}{}

\begin{itemize}
\tightlist
\item
  \hyperref[appendix-c99977d88b33]{Unsigned modular arithmetic} --- Modulo m retains the remainder class after division by m.
\item
  \hyperref[appendix-8f38927cfc54]{Bitwise NOT} --- Complement must specify width: eight-bit complement of \BookInlineCode{0x0F} is \BookInlineCode{0xF0}, while sixteen-bit complement is \BookInlineCode{0xFFF0}.
\end{itemize}

\subsection{Used by}\label{appendix-4fbc03d20829-used-by}

\hypertarget{appendix-4fbc03d20829-used-by}{}

\begin{itemize}
\item
  \hyperref[appendix-1d0dbe8924cc]{Fixed-point representation} --- uses this mechanism to define its representation, error or validation contract.
\item
  \hyperref[appendix-05f39b10364d]{Minimum signed value} --- uses two's complement as part of its representation or reasoning.
\item
  \hyperref[appendix-8d687c9c08ab]{Carry, overflow, negative, and zero flags} --- uses two's complement as part of its representation or reasoning.
\item
  \hyperref[appendix-1c7df340573e]{Subtraction through addition} --- uses two's complement as part of its representation or reasoning.
\item
  \hyperref[appendix-2ed044358b82]{Arithmetic policies} --- uses two's complement as part of its representation or reasoning.
\end{itemize}

\subsection{Recall and next depth}\label{appendix-4fbc03d20829-recall-and-next-depth}

\hypertarget{appendix-4fbc03d20829-recall-and-next-depth}{}

Explain two's complement without the label, work the example above, then state the assumption whose removal would change the result.

\textbf{Pending:} signed shifts, instruction semantics, overflow intrinsics, and wider arithmetic.

\subsection{Source and navigation}\label{appendix-4fbc03d20829-source-and-navigation}

\hypertarget{appendix-4fbc03d20829-source-and-navigation}{}

\hyperref[01-foundations-i-ch001-from-a-signal-to-an-interpretation]{C001 teaching chapter} · \hyperref[01-foundations-i-ch001-arithmetic-becomes-a-policy]{C002 development} · \href{../01-Bits-Programs-and-Reliable-Systems-P03/main.pdf\#nameddest=appendix-fd090d982a15}{Part definitions and mechanisms (Bits, Programs, and Reliable Systems, Part 3, Reference notes: Foundations)} · \href{../07-Vision-XR-and-Playable-Worlds-P01/main.pdf\#nameddest=appendix-4ad47bdd269a}{Dependency Map (Vision, XR, and Playable Worlds, Part 1, Reference notes: Dependency Map)}
